\documentclass[12pt,a4paper]{article}
\usepackage[T1]{fontenc}
\usepackage[utf8]{inputenc}
\usepackage[polish]{babel}
\usepackage{geometry}
\usepackage{booktabs}
\usepackage{longtable}
\usepackage{array}
\usepackage{amsmath}
\usepackage{graphicx}
\usepackage{hyperref}
\geometry{margin=2.5cm}

\title{Projekt 2: STRIPS\\Sprawozdanie z rozwiązania}
\author{
Autorzy:\\[0.5em]
Marcin Wolder\\
Stanisław Wojtas\\
}
\date{\today}

\begin{document}
\maketitle

\newpage

\section{Zakres pracy}
Niniejsze sprawozdanie obejmuje pełny zakres wymagań przewidzianych dla poziomów 4, 6 oraz 8 punktów:
\begin{itemize}
    \item definicję dziedziny STRIPS oraz 3 problemów w przestrzeni co najmniej 50 stanów,
    \item rozwiązanie metodą \texttt{forward planning} i zapis ciągów akcji,
    \item zaproponowanie heurystyki, uzasadnienie oraz porównanie czasów,
    \item zdefiniowanie podcelów (po co najmniej 2 na problem) i ponowne rozwiązanie z heurystyką i bez,
    \item trzy dodatkowe problemy z podcelami, w których plan wymaga co najmniej 20 akcji.
\end{itemize}

\section{Dziedzina i problemy bazowe}
W analizie przyjęto dziedzinę \texttt{parking}, w której występuje jedno puste miejsce \texttt{X}, a akcja elementarna polega na przemieszczeniu pojedynczego samochodu na pozycję pustą.
Dla 6 pozycji liczba stanów wynosi:
\[
6! = 720 > 50
\]
Źródło przykładowej domeny i instancji: 
\url{https://gist.github.com/primaryobjects/5ddc68b094d5affc492d2072e85be05b}

Problemy bazowe:
\begin{itemize}
    \item Problem A: start \texttt{\{1 2 3 X 4 5\}}, cel \texttt{\{5 1 X 3 2 4\}}, minimalna długość 6.
    \item Problem B: start \texttt{\{1 2 3 X 4 5\}}, cel \texttt{\{1 2 4 5 3 X\}}, minimalna długość 4.
    \item Problem C: start \texttt{\{1 2 3 X 4 5\}}, cel \texttt{\{1 3 2 4 5 X\}}, minimalna długość 5.
\end{itemize}

\section{Planowanie progresywne}
Do rozwiązania problemów zastosowano planowanie progresywne dostępne w bibliotece AIPython.

\begin{table}[h]
\centering
\resizebox{\textwidth}{!}{%
\begin{tabular}{lrrrr}
\toprule
Problem & Czas [s] & Rozszerzone ścieżki & Długość planu & Plan znaleziony \\
\midrule
A & 0.100878 & 616 & 6 & tak \\
B & 0.019958 & 203 & 4 & tak \\
C & 0.065413 & 459 & 5 & tak \\
\bottomrule
\end{tabular}
}
\caption{Wyniki wyszukiwania bez heurystyki.}
\end{table}

Uzyskane ciągi akcji (bez heurystyki):
\begin{itemize}
    \item Problem A:
    \begin{enumerate}
        \item \texttt{move\_c3\_s3\_s4}
        \item \texttt{move\_c5\_s6\_s3}
        \item \texttt{move\_c4\_s5\_s6}
        \item \texttt{move\_c2\_s2\_s5}
        \item \texttt{move\_c1\_s1\_s2}
        \item \texttt{move\_c5\_s3\_s1}
    \end{enumerate}
    \item Problem B:
    \begin{enumerate}
        \item \texttt{move\_c3\_s3\_s4}
        \item \texttt{move\_c4\_s5\_s3}
        \item \texttt{move\_c3\_s4\_s5}
        \item \texttt{move\_c5\_s6\_s4}
    \end{enumerate}
    \item Problem C:
    \begin{enumerate}
        \item \texttt{move\_c4\_s5\_s4}
        \item \texttt{move\_c3\_s3\_s5}
        \item \texttt{move\_c2\_s2\_s3}
        \item \texttt{move\_c3\_s5\_s2}
        \item \texttt{move\_c5\_s6\_s5}
    \end{enumerate}
\end{itemize}

\section{Heurystyka}
Zaproponowano heurystykę postaci:
\[
h(s)=\text{liczba samochodów na niewłaściwych miejscach (bez }X\text{)}
\]
Ponieważ pojedyncza akcja modyfikuje położenie dokładnie jednego pojazdu, wartość ta stanowi dolne ograniczenie liczby kroków koniecznych do osiągnięcia celu.

\begin{table}[h]
\centering
\resizebox{\textwidth}{!}{%
\begin{tabular}{lrrrrrr}
\toprule
Problem & Czas bez h [s] & Czas z h [s] & Rozszerz. bez h & Rozszerz. z h & Przyspieszenie & Dł. planu z h \\
\midrule
A & 0.082553 & 0.000303 & 532 & 7 & 272.86x & 6 \\
B & 0.019016 & 0.000168 & 204 & 5 & 112.91x & 4 \\
C & 0.056221 & 0.000204 & 307 & 6 & 275.03x & 5 \\
\bottomrule
\end{tabular}
}
\caption{Porównanie rozwiązań bez heurystyki i z heurystyką.}
\end{table}

Plany uzyskane z heurystyką:
\begin{itemize}
    \item Problem A:
    \begin{enumerate}
        \item \texttt{move\_c3\_s3\_s4}
        \item \texttt{move\_c4\_s5\_s3}
        \item \texttt{move\_c2\_s2\_s5}
        \item \texttt{move\_c1\_s1\_s2}
        \item \texttt{move\_c5\_s6\_s1}
        \item \texttt{move\_c4\_s3\_s6}
    \end{enumerate}
    \item Problem B:
    \begin{enumerate}
        \item \texttt{move\_c5\_s6\_s4}
        \item \texttt{move\_c3\_s3\_s6}
        \item \texttt{move\_c4\_s5\_s3}
        \item \texttt{move\_c3\_s6\_s5}
    \end{enumerate}
    \item Problem C:
    \begin{enumerate}
        \item \texttt{move\_c4\_s5\_s4}
        \item \texttt{move\_c5\_s6\_s5}
        \item \texttt{move\_c2\_s2\_s6}
        \item \texttt{move\_c3\_s3\_s2}
        \item \texttt{move\_c2\_s6\_s3}
    \end{enumerate}
\end{itemize}

\section{Planowanie z podcelami}
Dla każdego problemu zdefiniowano dwa podcele, a proces planowania realizowano etapowo:
\texttt{start \textrightarrow{} podcel1 \textrightarrow{} podcel2 \textrightarrow{} cel}.

Podcele:
\begin{itemize}
    \item A: \texttt{('c1','c2','c4','c3','X','c5')} oraz \texttt{('X','c1','c4','c3','c2','c5')}
    \item B: \texttt{('c1','c2','c3','c5','c4','X')} oraz \texttt{('c1','c2','c4','c5','X','c3')}
    \item C: \texttt{('c1','c2','c3','c4','c5','X')} oraz \texttt{('c1','c3','X','c4','c5','c2')}
\end{itemize}

\begin{table}[h]
\centering
\resizebox{\textwidth}{!}{%
\begin{tabular}{lrrrrrr}
\toprule
Problem & Czas bez h [s] & Czas z h [s] & Rozszerz. bez h & Rozszerz. z h & Przyspieszenie & Dł. planu z h \\
\midrule
A & 0.001306 & 0.000231 & 35 & 9 & 5.65x & 6 \\
B & 0.001014 & 0.000160 & 28 & 7 & 6.34x & 4 \\
C & 0.000701 & 0.000186 & 22 & 8 & 3.76x & 5 \\
\bottomrule
\end{tabular}
}
\caption{Wyniki dla planowania z podcelami (z i bez heurystyki).}
\end{table}

Uzyskane plany z podcelami i heurystyką są poprawne oraz zachowują długości planów z części bazowej (A: 6, B: 4, C: 5).

\subsection{Różnice czasów i przyspieszenia między podejściami}
W celu bezpośredniego porównania wariantów policzono:
\[
\Delta t = t_{\text{wolniejsze}} - t_{\text{szybsze}}, \qquad
\text{przyspieszenie} = \frac{t_{\text{wolniejsze}}}{t_{\text{szybsze}}}
\]

\begin{table}[h]
\centering
\resizebox{\textwidth}{!}{%
\begin{tabular}{lrrrrrr}
\toprule
Problem & \multicolumn{2}{c}{Podcele vs zwykłe (bez h)} & \multicolumn{2}{c}{Zwykłe + h vs zwykłe} & \multicolumn{2}{c}{Podcele + h vs zwykłe + h} \\
\cmidrule(lr){2-3}\cmidrule(lr){4-5}\cmidrule(lr){6-7}
 & $\Delta t$ [s] & Przysp. & $\Delta t$ [s] & Przysp. & $\Delta t$ [s] & Przysp. \\
\midrule
A & 0.099572 & 77.24x & 0.082250 & 272.86x & 0.000072 & 1.31x \\
B & 0.018944 & 19.68x & 0.018848 & 112.91x & 0.000008 & 1.05x \\
C & 0.064712 & 93.31x & 0.056017 & 275.03x & 0.000018 & 1.10x \\
\bottomrule
\end{tabular}
}
\caption{Bezpośrednie porównanie różnic czasów i przyspieszeń między wymaganymi wariantami.}
\end{table}

\section{Problemy rozszerzone (co najmniej 20 akcji)}
Zdefiniowano trzy dodatkowe problemy (D, E, F). Dla wszystkich przypadków przyjęto ten sam stan początkowy oraz dwa podcele.

\begin{table}[h]
\centering
\resizebox{\textwidth}{!}{%
\begin{tabular}{lrrrrrr}
\toprule
Problem & Czas bez h [s] & Czas z h [s] & Rozszerz. bez h & Rozszerz. z h & Długość planu z h & Warunek $\geq 20$ \\
\midrule
D & 0.337536 & 0.003187 & 2103 & 75 & 21 & tak \\
E & 0.337454 & 0.003274 & 2124 & 75 & 21 & tak \\
F & 0.337802 & 0.003271 & 2125 & 76 & 21 & tak \\
\bottomrule
\end{tabular}
}
\caption{Wyniki dla dodatkowych problemów z podcelami.}
\end{table}

Dla każdego z problemów D/E/F plan z heurystyką ma strukturę segmentową \texttt{7 + 7 + 7 = 21} akcji, więc wymaganie minimalnej długości 20 akcji jest spełnione.

\section{Wnioski}
\begin{itemize}
    \item Wszystkie wymagania zadania (zakres 4/6/8 punktów) zostały zrealizowane.
    \item We wszystkich przypadkach znaleziono poprawne plany.
    \item Heurystyka istotnie zmniejsza czas i liczbę rozszerzeń (szczególnie bez podcelów).
    \item Podejście z podcelami porządkuje rozwiązanie i pozwala skutecznie realizować dłuższe plany (21 akcji) dla problemów rozszerzonych.
\end{itemize}

\end{document}
